\begin{figure}[H]
    \centering
    % TikZ is placed in an LTR coordinate system.
    % Individual Persian text elements are restored with \faText.
    \begin{tikzpicture}
        \begin{axis}[
            width=0.85\textwidth,
            height=8.5cm,
            % Axis ranges
            ymin=0,
            ymax=88,
            xmin=0.15,
            xmax=14.85,
            % Background and axis lines
            axis background/.style={
                fill=Canvas
            },
            axis x line*=bottom,
            axis y line*=left,
            axis line style={
                draw=Ink,
                line width=0.75pt
            },
            tick style={
                draw=InkMuted,
                line width=0.55pt
            },
            % Y-axis
            ytick={0,20,40,60,80},
            yticklabels={
                \faText{۰},
                \faText{۲۰},
                \faText{۴۰},
                \faText{۶۰},
                \faText{۸۰}
            },
            ylabel={\faText{درصد فراوانی}},
            ylabel style={
                font=\bfseries\small,
                text=Ink,
                yshift=-2pt
            },
            yticklabel style={
                font=\small,
                text=InkMuted
            },
            % Grid
            xmajorgrids=false,
            ymajorgrids=true,
            grid style={
                draw=Border,
                line width=0.55pt
            },
            % One tick beneath every individual bar
            xtick={
                1,2,
                4,5,
                7,8,
                10,11,
                13,14
            },
            xticklabels={
                \faText{مرد},
                \faText{زن},
                \faText{مجرد},
                \faText{متاهل},
                \faText{شاغل},
                \faText{غیر شاغل},
                \faText{پایه},
                \faText{دانشگاهی},
                \faText{سطح کافی},
                \faText{سطح ناکافی}
            },
            xticklabel style={
                font=\footnotesize,
                text=Ink,
                rotate=52,
                anchor=east,
                align=right,
                yshift=-2pt
            },
            tick align=outside,
            % Bars and exact percentage labels
            bar width=13pt,
            clip=false,
            point meta=explicit symbolic,
            nodes near coords={\pgfplotspointmeta},
            every node near coord/.append style={
                font=\bfseries\scriptsize,
                text=Ink,
                anchor=south,
                yshift=2pt,
                inner sep=0.5pt
            }
        ]

        %--------------------------------------------------------------
        % First category of every independent variable
        %--------------------------------------------------------------
        \addplot+[
            ybar,
            mark=none,
            draw=PrimaryPurple,
            fill=PrimaryPurple
        ] coordinates {
            (1,56.2) [{۵۶٫۲٪}]
            (4,23.3) [{۲۳٫۳٪}]
            (7,56)   [{۵۶٪}]
            (10,39.5) [{۳۹٫۵٪}]
            (13,59.1) [{۵۹٫۱٪}]
        };

        %--------------------------------------------------------------
        % Second category of every independent variable
        %--------------------------------------------------------------
        \addplot+[
            ybar,
            mark=none,
            draw=SecondaryTeal,
            fill=SecondaryTeal
        ] coordinates {
            (2,22.2) [{۲۲٫۲٪}]
            (5,51.7) [{۵۱٫۷٪}]
            (8,38.1) [{۳۸٫۱٪}]
            (11,71)  [{۷۱٪}]
            (14,35.7) [{۳۵٫۷٪}]
        };

        %--------------------------------------------------------------
        % Independent-variable names and p-values
        %--------------------------------------------------------------
        \node[
            font=\bfseries\scriptsize,
            text=Ink,
            align=center,
            text width=1.9cm
        ]
        at (axis cs:1.5,-27) {
            \faText{جنسیت}\\[-1pt]
            \normalfont {$p<0.001$}
        };

        \node[
            font=\bfseries\scriptsize,
            text=Ink,
            align=center,
            text width=1.9cm
        ]
        at (axis cs:4.5,-27) {
            \faText{تاهل}\\[-1pt]
            \normalfont {$p=0.005$}
        };

        \node[
            font=\bfseries\scriptsize,
            text=Ink,
            align=center,
            text width=1.9cm
        ]
        at (axis cs:7.5,-27) {
            \faText{اشتغال}\\[-1pt]
            \normalfont {$p=0.029$}
        };

        \node[
            font=\bfseries\scriptsize,
            text=Ink,
            align=center,
            text width=1.9cm
        ]
        at (axis cs:10.5,-27) {
            \faText{تحصیلات}\\[-1pt]
            \normalfont {$p=$0.003}
        };

        \node[
            font=\bfseries\scriptsize,
            text=Ink,
            align=center,
            text width=1.9cm
        ]
        at (axis cs:13.5,-27) {
            \faText{سواد سلامت}\\[-1pt]
            \normalfont {$p=0.007$}
        };

        \end{axis}
    \end{tikzpicture}

    \caption{درصد فراوانی سلامت روانشناختی}
    \label{fig:psychological-health-grouped-bar-chart}
\end{figure}